%%%%%%%%%%%%%%%%%%%%%%%%%%%%%%%%%%%%%%%%%%%%%%%%%%%%%%%%%%%%
%% CHAPTER 11 -- Trust, Governance, and Who Is Responsible
%%%%%%%%%%%%%%%%%%%%%%%%%%%%%%%%%%%%%%%%%%%%%%%%%%%%%%%%%%%%%%%%%%

\chapter{Trust, Governance, and Who Is Responsible}
\label{chap:governance}
\label{ntg:ch11}

\scenelabel
\begin{scenestrip}
\sceneitem{The problem}
\sceneitem{How we got here}
\sceneitem{Where we stand}
\sceneitem{What goes wrong}
\end{scenestrip}

\paragraph{The problem.} Suppose you have built and deployed everything in the previous
ten chapters. Your model is pretrained, aligned, retrieval-
grounded, tool-equipped, and evaluated. It runs in production
and is being used by your customers, your employees, or your
students.

One morning a user takes a screenshot of the system giving a
confidently wrong piece of medical advice and posts it on
social media. The screenshot goes mildly viral. The next
morning, a regulator emails you with questions. The morning
after that, your insurance carrier asks for documentation of
your safety controls. The morning after that, a journalist
calls.

The technical question --- did the model produce the wrong
output? --- has a definite answer. The governance question ---
who is responsible, what process should have caught it, what
control will catch the next one, and how do you prove this to
a third party --- has a different shape. It is the question
this chapter is about. The good news, such as it is, is that
the discipline of answering it is older than the model by
half a century. The discipline is borrowed from aviation,
pharmaceuticals, financial services, and every other field
that has had to deploy systems whose failures matter. The
techniques have been adapted to AI. They are imperfect. They
are also the best tools we have.

\paragraph{How we got here.} Safety-critical engineering has been a serious discipline
since at least the 1970s. Nancy Leveson at MIT, in a series
of books and papers culminating in \emph{Engineering a Safer
World} in 2011, showed that accidents in complex systems are
typically not the result of a single component failure but of
the way the whole system is designed and operated. A door
falls off a plane not because one rivet was wrong but because
a series of small problems --- a documentation error, an
inspection oversight, a training gap, a procurement
shortcut --- aligned. The fix is not to find the one bad
rivet. The fix is to redesign the system so that small
problems do not align.

Charles Perrow's 1984 book \emph{Normal Accidents} extended
the argument: some systems are so tightly coupled that
accidents are statistically inevitable. The response is not
to wish the accidents away but to design systems whose
accidents are recoverable rather than catastrophic. James
Reason's \emph{Human Error} of 1990 gave us the Swiss-cheese
model of defence in depth: no single safety control is asked
to be perfect, but the alignment of holes across overlapping
controls is made improbable. Nuclear power plants, commercial
aviation, and pharmaceutical manufacturing all run on
versions of this insight. None of this work was about AI. All
of it applies to AI.

The translation to AI systems began with the model-card and
datasheet movement in the late 2010s. Margaret Mitchell at
Google, in a 2019 paper, proposed that every released model
should come with a structured document --- a ``model card''
--- describing what the model was trained on, what it was
evaluated against, what its known failure modes were, and what
uses it was and was not designed for. The same period saw the
``datasheets for datasets'' movement, which proposed similar
documentation for the training corpora themselves. Both
practices have been widely adopted, with varying levels of
rigour, by the major laboratories.

The regulatory translation began in the early 2020s. The
United States National Institute of Standards and Technology
(NIST) published the AI Risk Management Framework in January
2023, and a generative-AI profile of it the following year.
The framework gave US federal agencies and the contractors
that serve them a structured approach to AI risk. The
European Union's AI Act, agreed in 2024 and entering force in
stages, classified AI systems by risk tier and imposed
specific obligations on the high-risk and general-purpose
tiers. The UK established the AI Safety Institute, the United
States the AI Safety Institute, Singapore an evaluation
framework, and several other jurisdictions specific
disclosures and audit requirements. ISO/IEC 42001, published
in late 2023, gave organisations the first management-systems
standard for AI --- audit-ready in the same shape as ISO 9001
for quality or ISO 27001 for information security.

The combination of these frameworks is, in regulated sectors
and in the European Union, now the operating floor for
serious AI deployment. In less-regulated jurisdictions they
remain best practice rather than a hard requirement, but the
direction of travel is convergent. An organisation that
builds an LLM product today and is not aware of these
frameworks is building a product that will not be deployable
in many of the markets it might want to enter.

\paragraph{Where we stand.} A modern AI governance stack treats the deployed system as a
\emph{socio-technical artefact}: a model, plus an operator,
plus the users, plus the data flowing in and out, plus the
regulatory context that frames the whole. The engineering
disciplines map onto familiar shapes.

There is a \emph{hazard analysis}: a structured enumeration of
what can go wrong in the deployment, ranked by severity and
likelihood. The hazards include the obvious (the model
produces a wrong answer) and the less obvious (the model is
used to manipulate a vulnerable population, the model's
outputs are repurposed to build a system the operator did not
anticipate, the model's behaviour drifts as the underlying
vendor model is silently updated).

There is a set of \emph{controls}, layered for defence in
depth. Input filters that catch obvious abuse before the
model sees it. Retrieval grounding that gives the model
specific documents to cite. Alignment training that shapes
the model's default behaviour. Output classifiers that catch
the model's lapses before they reach the user. Human-in-the-
loop review for high-stakes outputs. Monitoring that flags
unusual patterns. Rollback procedures for when something has
gone wrong. None of the controls is asked to be perfect. The
combination is asked to be reliable.

There is an \emph{assurance case}: a structured argument,
with explicit claims, supporting arguments, and evidence,
that the controls are adequate for the deployment context.
The assurance case is read by internal review, by external
auditors, by regulators when applicable, and by the operator's
own future selves when something goes wrong and they need to
reconstruct what was claimed at deployment time. The
discipline of writing assurance cases is borrowed from the
aviation and nuclear industries; it is now being adapted to
AI.

There is \emph{documentation} --- model cards, datasheets,
algorithmic impact assessments, change logs, incident
reports --- that lets a third party reconstruct what was
built and why. This is the most boring part of the work and
the most consequential. A regulator who arrives at your door
is not asking for the model's source code. They are asking
for the documentation. If the documentation is good and
honest, the visit is a conversation. If it is bad or
dishonest, it is the beginning of a much longer process.

And there is \emph{monitoring}, because no static assessment
survives contact with the production traffic that follows it.
The behaviour of the deployed system is continuously
measured. Anomalies are surfaced. Regressions are caught.
The monitoring is the difference between a system that is
known to be safe and a system that is hoped to be safe.

\paragraph{What goes wrong.} The first failure of governance is to treat it as paperwork.
A model card that no one updates after the first deployment
is worse than no model card, because it gives false
reassurance. An assurance case written by the engineering
team to satisfy the legal team and never read by either is
worse than no assurance case, for the same reason. The
documentation is meant to be the externally legible artefact
of an actually-functioning process. Without the process, the
documentation is theatre.

The second failure is to trust a single control. Alignment
training alone is not enough; we have seen in Chapter~\ref{chap:alignment} that
it can be bypassed by jailbreaks. Output filtering alone is
not enough; clever attackers will find ways past the filter.
Retrieval grounding alone is not enough; we have seen in
Chapter~\ref{chap:systems} that it can be misled or injected against. The
Swiss-cheese model exists because every control has holes;
the discipline is to ensure the holes do not align.

The third failure is the moving target. A system that was
safe in October may not be safe in March, because the world
around it has changed even if the weights have not. New
jailbreak techniques have been published. New tool
integrations have been added. New user populations have
started using the system. The vendor has silently updated the
underlying model. The training data has aged. The regulatory
landscape has shifted. The governance work is not a one-time
event before deployment; it is an ongoing operation.

The fourth failure, and the one that is most distinctive of
this technology, is the \emph{honest limit}. There are
residual risk classes that alignment training and evaluation
alone do not fully eliminate. The model has no intrinsic
verifier of truth. The training data contains biases that the
training procedure does not penalise. The deployment
environment will produce inputs the evaluation did not
anticipate. An engineer who pretends otherwise is not being
useful. An engineer who acknowledges the limits and designs
the system to fail gracefully is being honest about the
problem. The frameworks and assurance disciplines in this
chapter exist to support that honesty.

\section{The idea, in plain language}

\subsection*{The aviation analogy and what it teaches}

The most useful single analogy for AI governance is
commercial aviation. Aviation was, in the early decades of
its history, dangerous. Aircraft fell out of the sky.
Passengers died. The industry's response, over decades, was
to build a culture and a regulatory apparatus that has made
commercial flight one of the safest forms of mass transit ever
devised. The mechanisms by which it did so are worth
understanding because they are roughly the mechanisms now
being adapted to AI.

The first mechanism is mandatory incident reporting. Every
near-miss, every unusual event, every minor incident is
reported to a central authority that aggregates the reports,
identifies patterns, and issues advisories. The reporting is
non-punitive within reasonable limits, because the
information is more valuable than the blame. AI is in the
early stages of building a comparable culture; the
\emph{adverse event} reporting mechanisms in the EU AI Act
and in some of the recent FDA guidance for AI-enabled
medical devices are the seeds of one.

The second mechanism is independent certification. An
aircraft cannot be sold for commercial use without being
certified by a regulatory authority that is institutionally
independent of the manufacturer. AI is, at the time of
writing, only beginning to have analogues of this. The EU AI
Act establishes a conformity-assessment regime for high-risk
AI systems. The various AI safety institutes around the world
are developing evaluation methodologies that may, over time,
become certification standards. The honest summary is that
the discipline is in its infancy, and the gap between what
the law currently requires and what would be sufficient is
wide.

The third mechanism is recurrent training and proficiency
review for the operators. A pilot does not, having qualified
once, fly forever; the qualifications are renewed regularly.
The analogous discipline for AI is the requirement that
operators of high-risk systems demonstrate ongoing
competence. This is mostly aspirational at this writing.

The fourth mechanism is the assumption that failures will
happen and the consequent design of systems to recover from
them. A commercial aircraft has redundant systems for every
critical function. When a system fails, the redundancy keeps
the aircraft flying long enough for the crew to land. The
analogous discipline for AI is the layered defence we have
already discussed: no single control is trusted, and the
system is designed to recover when one of them fails.

\subsection*{Who is responsible, and for what}

The legal and regulatory question of who is responsible when
an LLM produces harm is, at this writing, being worked out in
real time across jurisdictions. The general shape of the
emerging consensus is that responsibility is distributed
across several parties.

The \emph{provider} of the underlying model --- OpenAI,
Anthropic, Google, Meta, the labs that release open models ---
has obligations around how the model was trained, what it was
trained on, what its known capabilities and limits are, and
how it has been evaluated. These obligations are increasingly
codified in the EU AI Act and in vendor-specific terms of
service that allocate downstream responsibility.

The \emph{deployer} of the model --- the company that builds
a product on top of the model and offers it to users --- has
obligations around the appropriateness of the deployment, the
controls in place, the monitoring, and the response to
incidents. The deployer is, in most jurisdictions, the party
that will hear from regulators when something goes wrong,
even when the underlying issue is in the provider's model.

The \emph{user} of the deployed system --- the individual or
organisation that asks the model to do things --- has
obligations around what they ask of the system. A user who
asks the system to do something the deployer's terms of
service prohibit is not, in general, blameless when the
system complies.

There are also third parties who can be affected by the
system's outputs without ever having asked it to do anything.
A patient whose medical record was processed by an LLM-based
triage system. A defendant whose case was summarised by an
LLM for a judge. A job applicant whose r\'esum\'e was
screened by an LLM for an HR team. These third parties
typically have rights under the relevant regulations
(disclosure, contestability, data subject rights), and the
deployer has corresponding obligations.

\subsection*{The institutional architecture, briefly}

A non-technical reader who is going to engage with the
governance question --- as a policymaker, as an executive
overseeing an LLM deployment, as a journalist covering the
field --- benefits from knowing the rough institutional
architecture as it stands.

In the European Union, the AI Act establishes a
risk-tiered framework. Unacceptable-risk systems are
prohibited (social-scoring systems, certain real-time
biometric identification). High-risk systems (in education,
employment, critical infrastructure, certain
law-enforcement uses) face substantial conformity-assessment,
transparency, human-oversight, and post-market-monitoring
obligations. Limited-risk systems face transparency
obligations (such as disclosing that a user is interacting
with an AI). Minimal-risk systems face no specific
obligations beyond ordinary law.

In the United States, the regulatory landscape is more
fragmented. The NIST AI Risk Management Framework provides
voluntary guidance. Sector-specific regulators (the FDA for
medical devices, the FTC for consumer protection, the SEC
for financial services) apply their existing authorities to
AI within their sectors. State laws (notably California's)
are filling some of the gaps left by the absence of
comprehensive federal AI legislation. Executive orders and
agency guidance change with the administration.

In other jurisdictions, the picture is varied. The United
Kingdom is pursuing a sector-specific approach. China has
algorithm-registration and content-control requirements.
Singapore has a model AI governance framework. Brazil,
Canada, India, Japan, and others are at various stages of
their own legislative processes.

For an organisation deploying an LLM-based product
internationally, the governance work is to understand which
frameworks apply in each jurisdiction and to design the
deployment to satisfy the most stringent of the applicable
requirements. This is non-trivial. It is also the cost of
operating in this space.

\section{The engineering consequence}

\paragraph{Governance is engineering work, not legal work
alone.} The legal framework specifies what must be true. The
engineering work is what makes it true. A claim in a model
card that the model has been evaluated against a particular
hazard requires an evaluation harness that actually evaluates
against that hazard. A claim in an assurance case that an
output filter catches a particular class of misuse requires an
output filter that actually catches it. Treating governance as
a purely legal artefact, separate from the engineering, is
how organisations end up with documentation that does not
match reality.

\paragraph{Documentation is part of the product.} The model
card, the datasheet, the assurance case, the change log, the
incident report --- these are not external artefacts. They are
part of what is being shipped. A team that ships a model
without a model card has shipped an undocumented product. A
team that updates the model without updating the model card
has shipped a product whose documentation is now wrong. Both
are governance failures. Both are also operational risks: a
regulator who finds a discrepancy between the documentation
and the deployed behaviour will, reasonably, treat the
documentation as the binding statement.

\paragraph{Independent evaluation is increasingly standard.}
The frontier laboratories now subject their major models to
independent evaluation by external organisations before
release. The UK and US AI Safety Institutes, METR (formerly
ARC Evals), Apollo Research, and several others perform such
evaluations. The practice is not yet a regulatory requirement
in most jurisdictions, but it is rapidly becoming a market
expectation. An organisation deploying a model on behalf of
its users should ask whether the model has been independently
evaluated and, if so, by whom and against what.

\paragraph{Audit-readiness has a cost.} The engineering work
required to make a system audit-ready is substantial. The
costs include the documentation effort, the design of
evaluations that produce the right kinds of evidence, the
infrastructure for monitoring, and the organisational
processes for incident response. A small team building a
proof-of-concept can elide most of this. A team building a
production system in a regulated sector cannot. The transition
from one to the other is one of the more under-budgeted parts
of an LLM product's lifecycle.

\section{What goes wrong, and why}

\paragraph{Documentation drift.} The model card was accurate
on the day of deployment. Six months later it is not. The
model has been updated. The training data has been refreshed.
A new tool integration has been added. None of these were
reflected in the documentation. A regulator who asks for the
documentation receives a document that does not describe the
system in production. The fix is to treat the documentation as
versioned, like the code.

\paragraph{Single points of failure in the control stack.} A
system whose safety relies on a single control --- alignment,
or output filtering, or human review --- is a system whose
single point of failure is that control. The defence-in-depth
discipline calls for several controls, layered. The discipline
is sometimes skipped on cost grounds. The cost saving is
real; so is the resulting risk.

\paragraph{Treating evaluation as a release gate, not a
continuous practice.} A model that passes a battery of
evaluations at release time will, six months later, be facing
a different distribution of inputs and a different landscape
of attacks. The evaluations need to be re-run on a continuing
basis, against an evolving threat model. Static evaluation is
the equivalent of inspecting an aircraft once at the factory
and never again.

\paragraph{Vendor-side updates that the deployer does not
notice.} Vendors regularly update the models they serve under
the same API name. A deployer whose product is built on a
vendor's model can experience behaviour changes without
changing anything on their end. The fix is to pin to specific
model versions where the vendor allows, to monitor for
behavioural drift, and to test before any vendor-side update.
Not all vendors expose version pinning, and the deployers of
those vendors' models accept a real risk in doing so.

\paragraph{Automation complacency.} Commercial aviation, which
this chapter has used as its central analogy, achieved its
safety record in part by assuming that human operators will
become over-reliant on automation when systems are working
well. Pilots who have not manually flown an approach in months
are measurably slower to take control when the autopilot
fails. The corresponding risk in LLM deployments is that human
reviewers become progressively less critical as apparent
reliability improves. A review process that was rigorous when
the model was new will, over time, receive less scrutiny per
case. The error rate may not have changed; the human eye
trained to catch it has sharpened into complacency. The fix
in aviation is mandatory periodic practice under conditions
that do not allow automation to take over. The analogous
discipline for AI is to calibrate the review process to the
model's actual measured error rate, not to reviewers'
subjective confidence in a system they have grown used to.

\paragraph{The honest limits of evaluation.} The current
evaluation methodology is good at catching the failure modes
the evaluators thought to test for. It is, by its nature, less
good at catching failure modes nobody thought to test for. The
gap between ``evaluated'' and ``safe'' is real, and any
governance framework that pretends otherwise is overselling
itself. The honest position, for now, is that we evaluate
against the best understanding of the threat model we have,
that we update the threat model as new failures are
discovered, and that we accept some residual risk that no
amount of evaluation will reduce to zero.

\section*{Bridges}
\addcontentsline{toc}{section}{Bridges}

This chapter built on every preceding chapter, because
governance is the discipline of overseeing the system that
the previous chapters described how to build. Chapter~\ref{chap:agents}
returns to the developer-facing surface of these systems and
the specific attack patterns the agentic versions of them
expose. Chapter~\ref{chap:error-model} returns to the failure-mode vocabulary of
this chapter and reorganises it into a single ladder a
non-technical reader can use to diagnose particular outputs.
The next chapter --- \emph{How You Talk to a Model Today} ---
is the closer look at the developer-facing layer that ties
together all the engineering of the previous chapters into the
products you actually use.
